\documentclass[12pt,a4paper]{article}
\usepackage[utf8]{inputenc}
\usepackage[T1]{fontenc}
\usepackage[left=1.4cm,right=1.4cm,top=0.95cm,bottom=0.95cm]{geometry}
\usepackage[hidelinks]{hyperref}
\usepackage{enumitem}
\usepackage{titlesec}
\pagestyle{empty}
\setlength{\parindent}{0pt}
\titleformat{\section}{\large\bfseries}{}{0em}{}[\titlerule]
\titlespacing*{\section}{0pt}{3.5pt}{2pt}
\setlist[itemize,1]{leftmargin=0pt,label={},topsep=0.5pt,itemsep=1.5pt,parsep=0pt}
\setlist[itemize,2]{leftmargin=1.2em,label=\textbullet,topsep=0pt,itemsep=0pt,parsep=0pt}

\begin{document}

{\centering
{\LARGE\textbf{Nicolas Rojas Varela}}\\[0.5mm]
{\large\textbf{AI R\&D Engineer | Research Engineer}}\\[0.5mm]
Annecy, Auvergne-Rhône-Alpes, France\\
+33 7 45 61 10 20 \textbar{} \href{mailto:nicolasrojas1111@gmail.com}{nicolasrojas1111@gmail.com} \textbar{} \href{https://github.com/nirojasva}{GitHub} \textbar{} \href{https://nirojasva.github.io/cv/}{Online CV}\\
Authorized to work in France \textbar{} Available after PhD defense (March 2027)\par}
\vspace{1pt}

\section*{Profile}
Computer Science Researcher and R\&D Engineer specialized in machine learning applied to time series, data streams, and anomaly detection. Experience in algorithm design, experimental evaluation, and applied research in academic and industrial contexts. Industrial PhD (CIFRE) candidate at Université Clermont Auvergne (LIMOS), with a defense scheduled for March 2027 (\href{https://spectral-anomaly-detector.limos.fr/index_en.html}{Thesis Webpage}). Seeking an AI R\&D Engineer or Research Engineer position.

\section*{Research Expertise}
\textbf{Domains:} Anomaly detection, unsupervised learning, data streams, time series, spectral data, classification, interpretability\\
\textbf{Methods:} K-NN, Isolation Forest, autoencoders, model benchmarking, experimental reproducibility\\
\textbf{Bibliothèques Python:} scikit-learn, CapyMOA, PySAD, Aeon, sktime, PyTorch, TensorFlow\\
\textbf{Additional Tools:} SQL, R, Git, Linux, GCP, data pipelines, ETL, APIs

\section*{R\&D and Research Experience}
\begin{itemize}
\item \textbf{R\&D Engineer / Industrial PhD Researcher} \hfill March 2024 -- March 2027 (expected)\\
\textit{Université Clermont Auvergne - LIMOS (CIFRE)} \hfill Annecy and Aubière, France
\begin{itemize}
\item Design and evaluation of unsupervised methods for real-time spectral anomaly detection in an industrial environment (OnlineBootKNN model: 85\% detection rate and 0.5\% false positive rate).
\item Formulation of online learning approaches on spectral data streams, optimizing algorithmic efficiency and model interpretability (processing time under 1 s per spectrum).
\item Development of reproducible experimental pipelines in Python using CapyMOA, scikit-learn, and PySAD.
\end{itemize}

\item \textbf{Machine Learning Researcher} \hfill November 2023 -- March 2024\\
\textit{LIMOS} \hfill Aubière, France
\begin{itemize}
\item Development of time series classification and anomaly detection algorithms for data streams (interpretable shapelet-based RSAST method, 3 min 15 s execution time on time series of length 1024).
\item Design of comparative experimental protocols, model evaluation, and open-source integration using Aeon, scikit-learn, and sktime.
\end{itemize}

\item \textbf{Data Mining Intern} \hfill March 2023 -- September 2023\\
\textit{LIMOS} \hfill Aubière, France
\begin{itemize}
\item Interpretable time series analysis and classification using Python methods and libraries, notably Aeon, scikit-learn, and sktime.
\end{itemize}
\end{itemize}

\section*{Scientific Publications}
\begin{itemize}
\item \textbf{Rojas Varela, N.}, Ah-Pine, J., \& Mephu Nguifo, E. (2026). \textit{OnlineBootKNN: An Unsupervised Framework for Detecting Anomalies in Spectral Data Streams}. \textit{Proceedings of the AAAI Conference on Artificial Intelligence}, 40(18), 15797--15805. \href{https://doi.org/10.1609/aaai.v40i18.38611}{DOI} \textbar{} \href{https://github.com/nirojasva/spectral-benchmark}{Code}
\item \textbf{Rojas Varela, N.}, Ah-Pine, J., \& Mephu Nguifo, E. (2026). \textit{The BV4 Benchmark for Unsupervised Anomaly Detection in High-Dimensional Spectral Data Streams}. \textit{Advances in Neural Information Processing Systems (NeurIPS), Datasets and Benchmarks Track}. Accepted, to appear. \href{https://huggingface.co/datasets/nirojasva/spectral-benchmark-bv4}{Hugging Face}
\end{itemize}

\section*{Workshops and Presentations}
\begin{itemize}
\item Amroun, C., \textbf{Rojas Varela, N.}, \& Mephu Nguifo, E. (2026). \textit{Enhanced OBKNN: An Efficient Framework for Anomaly Detection in Spectral Data Streams}. \textit{40th Conference on Neural Information Processing Systems (NeurIPS 2026), Workshop on TS-LIMITS}. Accepted, to present. \href{https://hal.science/hal-05803665/}{HAL}.
\item \textbf{Rojas Varela, N.}, Ah-Pine, J., \& Mephu Nguifo, E. (2026). \textit{OnlineBootKNN: An Unsupervised Framework for Detecting Anomalies in Spectral Data Streams}. \textit{ECML PKDD 2026 (Nectar Track)}, oral presentation. Short version of the AAAI paper listed above.
\item \textbf{Rojas Varela, N.}, Ah-Pine, J., \& Mephu Nguifo, E. (2024). \textit{Study and design of interpretable models for detecting anomalies in spectral data streams}. \textit{ECML PKDD 2024 (PhD Forum Track)}.
\item \textbf{Rojas Varela, N.}, Mbouopda, M. F., \& Mephu Nguifo, E. (2023). \textit{RSAST: Sampling Shapelets for Interpretable Time Series Classification}. Presented at \textit{Conférence Francophone sur l'Apprentissage Automatique (CAP)}. \href{https://hal.science/hal-04311309/}{HAL} \textbar{} \href{https://www.aeon-toolkit.org/en/v1.1.0/api_reference/auto_generated/aeon.transformations.collection.shapelet_based.RSAST.html}{Code (Aeon)}
\end{itemize}

\section*{Education}
\begin{itemize}
\item \textbf{PhD in Computer Science} -- \textit{LIMOS, Université Clermont Auvergne} \hfill 2024 -- 2027
\item \textbf{Master International in Computer Science} -- \textit{ISIMA} \hfill 2022 -- 2023
\item \textbf{Specialization in Statistics} -- \textit{Universidad Nacional de Colombia} \hfill 2020 -- 2021
\item \textbf{BSc in Industrial Engineering} -- \textit{Universidad Nacional de Colombia} \hfill 2012 -- 2017
\end{itemize}

\section*{Additional Professional Experience}
\begin{itemize}
\item \textbf{Business Intelligence Analyst} \hfill October 2021 -- August 2022\\
\textit{Center Source, Inc.} \hfill Barranquilla, Colombia
\begin{itemize}
\item Design of ETL pipelines and data modeling (tracking key business customer indicators updated in real time).
\item Development of dashboards and analytics solutions using Python, R, Tableau, and Exasol.
\end{itemize}

\item \textbf{Data Analyst} \hfill August 2018 -- October 2021\\
\textit{Seguros Bolívar S.A.} \hfill Bogotá, Colombia
\begin{itemize}
\item Statistical modeling, predictive analytics, and large-scale data analysis (improving customer service levels by up to 80\%).
\item Use of Python, R, Excel, GCP, Dataiku, and Power BI to support operational activities.
\end{itemize}
\end{itemize}

\section*{Languages}
Spanish: Native \textbar{} English \textbar{} French

\end{document}